\section{Multilinearidade}

Em diante abreviamos \(\Pi V_j = V_1 \times V_2 \times \cdots \times V_k\) como o produto cartesiano dos \(\F\)-espaços vetoriais \(V_j\) com \(j\leq k \in \N\). 
%    Sejam \(\Pi V_j\) e \( W\) \(\F\)-espaço vetorias. Denotamos por \( \mathcal{F}\left(\Pi V_j, W\right)\) o conjunto de todas as funções \(f : \Pi V_j \to W \).    
%\end{definition}
\begin{definition}
    Sejam \(W\) \(\F\)-espaço. Uma função \(\phi \in \mathcal{F}( \Pi V_j, W)\) é $k$\emph{-linear} se é linear em cada entrada. Ou seja, \(\forall j_0\leq k\), \(\forall (v_j)_{j\neq j_0}\in \Pi V_j \) fixos, \(\forall v_{1,j_0}, v_{2,j_0} \in V_{j_0} \) e \( \forall \lambda \in \F\),   
    \[  \textstyle \phi_{j_0} (\cdots, v_{1,j_0} + \lambda v_{2,j_0}, \cdots ) = \phi ( \cdots, v_{1,j_0}, \cdots) + \lambda \phi(\cdots, v_{2,j_0}, \cdots ).\]  
\end{definition}

\vspace{-0.5cm}
\begin{definition}
    \centering \(\hom^k(\Pi V_j,W):= \{\phi \in \mathcal{F}(\Pi V_j, W): \phi \text{ é $k$-linear}\}\).
\end{definition}

\begin{note}
    Se \(\forall j\leq k\), \(V_j =V\), então \(\hom^k(\Pi V_j,W) = \hom^k(V,W)\). Se \(W = \F \), chamamos seus elementos de \emph{formas $k$-lineares} em \(V\). 
\end{note}

\alerta{
    \centering \(U \leq \Pi V_j\  \textcolor{red}{\nRightarrow} \ \phi(U) \leq W\)
    \tcblower
    \begin{example}
        Sejam \(\phi : \R^2 \times \R^2 \ni (u,v)  \mapsto u \cdot v^t \in M_2(\R)\). \comentario{Veja que é bilinear, \(\operatorname{Im}\phi = \{A \in M_2(\R) : \operatorname{pt} A = 1\}\). Ahí \( E_{11}+E_{22} = \id_2 \notin \operatorname{Im}\phi\) } 
    \end{example}
}

\theoremnum{9.1.3.}
\begin{theorem}
    \label{thm:9.1.3}
    \(\forall j\leq k\) seja \(\alpha_j = (v_{i,j})_{i \in I_j}\) base de \(V_j\), \(I = \Pi^k I_j\) e \((w_\textbf{i})_{\textbf{i} \in I}\) família de vetores em \(W\). Então, \(\exists ! \phi \in \hom^k(\Pi V_j, W):\forall \textbf{i} \in I\) temos \(\phi((v_{\textbf{i}, j})) = w_\textbf{i}\). 
\end{theorem}

\newcommand{\bi}{\textbf{i}}
\demo{Sejam \((v_{j}) \in \Pi V_j\) e \(\forall j \leq k\), \( (a_{i_{j},j}) = [v_{{j}}]_{\alpha_{j}}\) com \(i_j \leq m_j \in \mathbb{Z}_{\geq 0} \), supondo \(\phi\) seja $k$-linear e \(I_m:=\prod I_{m_j} \) mostra primeiro que,     
\[ \phi((v_{j})) = \sum_{\bi \in I_m} \left({\prod}^k a_{i_j,j}\right) \cdot  \phi((v_{\bi,j})).\]
Use aquela expressão pra replicar a demonstração do \(\blacksquare\) \ref{thm:6.1.6}. 
}%A unicidade segue da unicidade das coordenadas. Agora, com \(j_0\) fixo olha pra \(\phi((v_{1,j_0}) + \lambda(v_{2,j_0}))\).  } = \sum_{\bi \in I_m} \left(\prod a_{i_j,j}\right) \cdot  w_\bi
%\demo{A demostração é completamente análoga da feita pra \(\blacksquare\) \ref{thm:6.1.6}, mudando, claro, o não despreciavél negôcio da gestão dos índices}

\propositionnum{9.1.4.}
%, \(\textbf{v}_\textbf{i}:= (v_{\textbf{i},j})_{\textbf{i} \in I}\)
\begin{proposition}
    \label{prop:9.1.4}
    Sejam \(V_j\) e \(\alpha_j\) como no \(\blacksquare\) \ref{thm:9.1.3}, \(\beta = (w_s)_{s\in S}\) uma base de \(W\) e \(\forall (\textbf{i}, s) \in I \times S\), \(\phi_{\textbf{i}, s} \in \hom^k(\Pi V_j, W)\) o único tal que \(\forall \tilde{\textbf{i}} \in I, \phi_{\textbf{i},s}((v_{\tilde{\textbf{i}}, j })) = \delta_{\textbf{i}, \tilde{\textbf{i}}} \ w_s\). Então, \((\phi_{\textbf{i}, s})\) é l.i., mas ainda, se \(\forall j\leq k, \ \dim V_j < \infty \), a família é base do \(\hom^k(\Pi V_j, W)\). 
\end{proposition}

\comentario{
    A demonstração não suporta resumo, mas segue sendo o mesmo negócio de sempre, pega \(\Gamma = \{\gamma_1, \cdots, \gamma_m\} \subseteq I \times S\) finita e mostra que \(\sum^m a_l \phi_{\gamma_l} = 0 \) \(\leftrightharpoons\) \(\forall l \leq m, \ a_l = 0 \), olha se quiser \cite[Pág. 299]{MA719}
}

\begin{corollary}
    Se \(\forall j\leq k, \dim V_j<\infty\), então \(\dim \hom^k(\Pi V_j, W) = \dim W \cdot \prod \dim V_j\). \footnote{Se \(\exists j_0 \leq k: \dim V_j = \infty\), então \(\dim \hom^k(\Pi V_j, W) > \dim W \cdot \prod \dim V_j\). \comentario{Não é trivial}}
\end{corollary}
